\documentclass{ximera}
%nieuw 2 september 2026
\title{Samenvatting: het elektrisch veld}
\author{Daan Lipkens}

\begin{document}
\begin{abstract}
    Samenvatting van hoofdstuk 2: het elektrisch veld.
\end{abstract}

\begin{onlineOnly}%voor bladindeling
    \maketitle
\end{onlineOnly}

\newgeometry{left=3cm,bottom=0.1cm}
\pagestyle{empty}

\begin{conclusion}[title={Samenvatting}]\nl

\samenbox{Het elektrisch veld}{
    Een bronlading $Q_b$ creëert in de ruimte rondom zich een elektrisch veld $\vec{E}$. Een testlading $Q_t$ in dat veld ondervindt een kracht $\vec{F}_e = Q_t \cdot \vec{E}$. Het veld zelf is onafhankelijk van de aanwezigheid van een testlading.
    $$ \vec{E} \equiv \frac{\vec{F}_e}{Q_t} \qquad \text{eenheid: } \mathrm{N/C} $$
}

\samenbox{Veld van een puntlading (radiaal veld)}{
    $$ E = k_e \frac{|Q_b|}{r^2} $$
    \begin{itemize}
        \item Positieve bronlading $\rightarrow$ veld wijst radiaal \textbf{weg} van de lading.
        \item Negatieve bronlading $\rightarrow$ veld wijst radiaal \textbf{naar} de lading toe.
    \end{itemize}
}

\samenbox{Superpositiebeginsel}{
    Het totale veld van meerdere bronladingen is de \textbf{vectoriële som} van de afzonderlijke velden:
    $$ \vec{E}_{tot} = \vec{E}_1 + \vec{E}_2 + \vec{E}_3 + \dots $$
    Liggen alle bronladingen en het punt \textit{P} op één rechte lijn, dan volstaat een algebraïsche optelling.
}

\samenbox{Veldlijnen}{
    \begin{itemize}
        \item Raaklijnig aan $\vec{E}$ in elk punt; pijl geeft de zin aan.
        \item Dichtheid van de lijnen $\propto$ grootte van het veld.
        \item Vertrekken bij een positieve lading, eindigen bij een negatieve lading.
        \item Aantal lijnen $\propto$ grootte van de lading.
        \item Veldlijnen kunnen elkaar \textbf{nooit} snijden.
    \end{itemize}
    Een \textbf{homogeen veld} heeft overal dezelfde grootte, richting en zin (evenwijdige, gelijk verdeelde veldlijnen), zoals tussen twee geladen platen.
}

\samenbox{Materie in een elektrisch veld}{
    \begin{itemize}
        \item Een \textbf{isolator} polariseert lichtjes in een elektrisch veld.
        \item In een \textbf{geleider} in elektrostatisch evenwicht is het veld \textbf{binnenin} steeds \textbf{nul}, en staan de veldlijnen aan het oppervlak \textbf{loodrecht}.
        \item Een holle geleider schermt zijn binnenkant af van externe velden: \textbf{elektrische schermwerking} (kooi van Faraday). Toepassingen: coaxkabel, auto tijdens onweer, afscherming van gevoelige elektronica.
    \end{itemize}
}

\end{conclusion}

\restoregeometry
\pagestyle{plain}

\end{document}
